% Part: first-order-logic
% Chapter: models-theories
% Section: set-theory

\documentclass[../../../include/open-logic-section]{subfiles}

\begin{document}

\olfileid{fol}{mat}{set}

\olsection{نظریهٔ مجموعه‌ها}

تقریباً همهٔ ریاضیات را می‌توان در نظریهٔ مجموعه‌ها بسط داد. بسط
ریاضیات در این نظریه مستلزم چند کار است. نخست، به مجموعه‌ای از اصول
موضوع برای رابطهٔ~$\in$ نیاز داریم. چندین دستگاه متفاوت از اصول موضوع
پدید آمده‌اند که گاه خواص متعارضی را برای~$\in$ مقرر می‌کنند. در این
میان، دستگاه اصول موضوع موسوم به~$\Log{ZFC}$، یعنی نظریهٔ مجموعه‌های
زرملو--فرانکل همراه با اصل موضوع انتخاب، برجسته است: این دستگاه با
فاصلهٔ بسیار پرکاربردترین و پرمطالعه‌ترین دستگاه است، زیرا معلوم شده
اصول موضوع آن برای اثبات تقریباً همهٔ چیزهایی که ریاضی‌دانان انتظار
دارند بتوانند اثبات کنند کافی است. اما پیش از آنکه این مطلب را بتوان
نشان داد، نخست باید روشن کنیم که اساساً چگونه می‌توانیم همهٔ چیزهایی
را \emph{بیان} کنیم که ریاضی‌دانان می‌خواهند بیان کنند. برای آغاز، زبان
هیچ !!{constant} یا !!{function}ی ندارد؛ پس در نگاه نخست معلوم نیست
چگونه می‌توانیم دربارهٔ مجموعه‌های خاص (مانند~$\emptyset$ یا~$\Nat$)،
عملیات روی مجموعه‌ها (مانند~$X \cup Y$ و~$\Pow{X}$)، چه رسد به
ساخت‌هایی چون روابط و توابع که چیزهایی غیر از مجموعه‌ها را در بر
می‌گیرند، سخن بگوییم.

برای شروع، «عضوِ ... است» تنها رابطهٔ مورد نظر ما نیست؛ «زیرمجموعهٔ
... است» تقریباً به همان اندازه مهم به نظر می‌رسد. اما می‌توانیم
«زیرمجموعهٔ ... است» را بر حسب «عضوِ ... است» \emph{تعریف} کنیم. برای
این کار باید !!a{formula} چون~$!A(x, y)$ در زبان نظریهٔ مجموعه‌ها بیابیم
که زوج مجموعه‌های~$\tuple{X, Y}$ آن را ارضا کنند هرگاه و تنها هرگاه
$X \subseteq Y$. اما $X$ دقیقاً هنگامی زیرمجموعهٔ~$Y$ است که همهٔ
!!{element}s~$X$، !!{element}s~$Y$ نیز باشند. پس می‌توانیم~$\subseteq$
را با فرمول زیر تعریف کنیم:
\[
\lforall[z][(z \in x \lif z \in y)]
\]
اکنون هرگاه بخواهیم رابطهٔ~$\subseteq$ را در فرمولی به کار ببریم،
می‌توانیم به‌جای آن از این فرمول استفاده کنیم (با جایگزینی مناسب~$x$ و
$y$ و، در صورت لزوم، تغییر نام متغیر مقید~$z$). برای نمونه، گسترشی‌بودن
مجموعه‌ها به این معناست که اگر هر یک از دو مجموعهٔ~$x$ و~$y$ در دیگری
مندرج باشد، آن دو باید یک مجموعه باشند. این مطلب را می‌توان با
$\lforall[x][\lforall[y][((x \subseteq y \land y
    \subseteq x) \lif x = y)]]$ بیان کرد؛ یا، اگر~$\subseteq$ را با
تعریف بالا جایگزین کنیم، با فرمول زیر:
\[
\lforall[x][\lforall[y][((\lforall[z][(z \in x \lif z \in y)] \land
    \lforall[z][(z \in y \lif z \in x)]) \lif x = y)]].
\]
این در واقع یکی از اصول موضوع~$\Log{ZFC}$، یعنی «اصل گسترش»، است.

هیچ !!{constant}ی برای~$\emptyset$ وجود ندارد، اما می‌توانیم با
$\lnot \lexists[y][y \in x]$ بیان کنیم که «$x$ تهی است». در این صورت
عبارت «$\emptyset$ وجود دارد» به !!{sentence}
~$\lexists[x][\lnot \lexists[y][y \in x]]$ بدل می‌شود. این نیز یکی دیگر
از اصول موضوع~$\Log{ZFC}$ است. (توجه کنید که اصل گسترش مستلزم آن است که
تنها یک مجموعهٔ تهی وجود داشته باشد.) هرگاه بخواهیم در زبان نظریهٔ
مجموعه‌ها دربارهٔ~$\emptyset$ سخن بگوییم، باید آن را به صورت «مجموعه‌ای
وجود دارد که تهی است و~\dots» بنویسیم. برای نمونه، برای بیان اینکه
$\emptyset$ زیرمجموعهٔ هر مجموعه است می‌توانیم بنویسیم
\[
\lexists[x][(\lnot\lexists[y][y \in x] \land \lforall[z][x \subseteq
    z])]
\]
که البته در آن باید $x \subseteq z$ را نیز به نوبهٔ خود با تعریفش
جایگزین کنیم.

برای سخن‌گفتن از عملیات روی مجموعه‌ها، مانند~$X \cup Y$ و~$\Pow{X}$،
باید شگرد مشابهی به کار ببریم. زبان نظریهٔ مجموعه‌ها هیچ نماد تابعی
ندارد، اما می‌توانیم روابط تابعی~$X \cup Y = Z$ و~$\Pow{X} = Y$ را با
فرمول‌های زیر بیان کنیم:
\begin{align*}
& \lforall[u][((u \in x \lor u \in y) \liff u \in z)]\\
& \lforall[u][(u \subseteq x \liff u \in y )]
\end{align*}
زیرا !!{element}s~$X \cup Y$ دقیقاً مجموعه‌هایی هستند که یا
!!{element}s~$X$ یا !!{element}s~$Y$ هستند، و !!{element}s~$\Pow{X}$
دقیقاً زیرمجموعه‌های~$X$ هستند. بااین‌حال، این کار به ما اجازه نمی‌دهد
$x \cup y$ یا~$\Pow{x}$ را چنان به کار ببریم که گویی ترم‌اند؛ تنها
می‌توانیم کل !!{formula}sی را به کار ببریم که روابط~$X \cup Y = Z$ و
$\Pow{X} = Y$ را تعریف می‌کنند. در واقع، هنوز نمی‌دانیم این روابط اصلاً
ارضا می‌شوند یا نه؛ یعنی نمی‌دانیم اجتماع‌ها و مجموعه‌های توانی همواره
وجود دارند. برای نمونه، !!{sentence}
$\lforall[x][\lexists[y][\Pow{x} = y]]$ یکی دیگر از اصول موضوع
$\Log{ZFC}$ است (اصل موضوع مجموعهٔ توانی).

حال سخن‌گفتن از زوج‌های مرتب یا توابع چه می‌شود؟ در اینجا باید توضیح
دهیم که چگونه می‌توان زوج‌های مرتب و توابع را انواع خاصی از مجموعه‌ها
دانست. یکی از راه‌های تعریف زوج مرتب~$\tuple{x, y}$ این است که آن را
مجموعهٔ~$\{\{x\}, \{x, y\}\}$ بگیریم. اما مانند قبل نمی‌توانیم
!!a{function} معرفی کنیم که این مجموعه را نام‌گذاری کند؛ فقط می‌توانیم
رابطهٔ~$\tuple{x, y} = z$، یعنی~$\{\{x\}, \{x, y\}\} = z$، را تعریف
کنیم:
\[
\lforall[u][(u \in z \liff (\lforall[v][(v \in u \liff v = x)] \lor
  \lforall[v][(v \in u \liff (v = x \lor v = y))]))]
\]
این فرمول می‌گوید !!{element}s~$u$ از~$z$ دقیقاً مجموعه‌هایی هستند که
یا تنها~$x$ را به‌عنوان !!{element} خود دارند یا فقط~$x$ و~$y$ را
به‌عنوان !!{element}s خود دارند (به بیان دیگر، مجموعه‌هایی که یا با
$\{x\}$ یکسان‌اند یا با~$\{x, y\}$). پس از داشتن این تعریف، می‌توانیم
چیزهای بیشتری بگوییم؛ برای نمونه، اینکه $X \times Y = Z$:
\[
\lforall[z][(z \in Z \liff \lexists[x][\lexists[y][(x \in
        X \land y \in Y \land \tuple{x, y} = z)]])]
\]

می‌توان تابع~$f \colon X \to Y$ را رابطهٔ~$f(x) = y$، یعنی مجموعهٔ
زوج‌های~$\Setabs{\tuple{x,y}}{f(x) = y}$، دانست. سپس می‌توانیم بگوییم
مجموعهٔ~$f$ تابعی از~$X$ به~$Y$ است اگر (الف) رابطه‌ای باشد که
$\subseteq X \times Y$، (ب) تام باشد، یعنی برای هر~$x \in X$ یک
$y \in Y$ وجود داشته باشد چنان‌که~$\tuple{x, y} \in f$، و (پ)
تک‌مقداری باشد، یعنی هرگاه~$\tuple{x, y}, \tuple{x, y'} \in f$، آنگاه
$y = y'$ (زیرا مقادیر تابع باید یکتا باشند). بنابراین «$f$ تابعی از
$X$ به~$Y$ است» را می‌توان چنین نوشت:
\begin{align*}
 \lforall[u][(u \in f \lif {}] & \lexists[x][\lexists[y][(x \in X \land y \in
      Y \land \tuple{x, y} = u)]]) \land {}\\
 \lforall[x][(x \in X \lif {}] &
      (\lexists[y][(y \in Y \land \mathrm{maps}(f, x, y))] \land {}\\
& (\lforall[y][\lforall[y'][((\mathrm{maps}(f, x, y) \land
    \mathrm{maps}(f, x, y')) \lif y = y')]]))
\end{align*}
که در آن $\mathrm{maps}(f, x, y)$ صورت کوتاه
$\lexists[v][(v \in f \land \tuple{x, y} = v)]$ است (این !!{formula}
عبارت «$f(x) = y$» را بیان می‌کند).

اکنون بیان اینکه $f\colon X \to Y$، برای نمونه، !!{injective} است نیز
دشوار نیست:
\begin{multline*}
 f \colon X \to Y \land \lforall[x][\lforall[x'][((x \in X \land x' \in
     X \land {}]] \\
 \lexists[y][(\mathrm{maps}(f, x, y) \land \mathrm{maps}(f,
        x', y))]) \lif x = x')
\end{multline*}
تابع~$f\colon X \to Y$ !!{injective} است هرگاه و تنها هرگاه، اگر~$f$ دو
عضو~$x, x' \in X$ را به یک~$y$ بفرستد، آنگاه $x = x'$. اگر این فرمول را
با~$\mathrm{inj}(f, X, Y)$ کوتاه کنیم، همین حالا نیز می‌توانیم مطلبی
به پیچیدگی قضیهٔ کانتور را در زبان نظریهٔ مجموعه‌ها بیان کنیم: هیچ
!!{injective} تابعی از~$\Pow{X}$ به~$X$ وجود ندارد:
\[
\lforall[X][\lforall[Y][(\Pow{X} = Y \lif
    \lnot\lexists[f][\mathrm{inj}(f, Y, X)])]]
\]

ممکن است گمان کنیم نظریهٔ مجموعه‌ها به اصل موضوع دیگری نیاز دارد که
وجود مجموعه‌ای را برای هر ویژگی تعریف‌کننده تضمین کند. اگر~$!A(x)$
فرمولی از نظریهٔ مجموعه‌ها با متغیر آزاد~$x$ باشد، می‌توانیم
!!{sentence} زیر را در نظر بگیریم:
\[
\lexists[y][\lforall[x][(x \in y \liff !A(x))]].
\]
این !!{sentence} می‌گوید مجموعه‌ای چون~$y$ وجود دارد که !!{element}s آن
همه و تنها همان~$x$هایی هستند که~$!A(x)$ را ارضا می‌کنند. این طرح‌واره
«اصل مجموعه‌سازی» نامیده می‌شود. بسیار سودمند به نظر می‌رسد؛ اما
متأسفانه ناسازگار است. بگذارید~$!A(x) \ident \lnot x \in x$؛ آنگاه اصل
مجموعه‌سازی می‌گوید
\[
\lexists[y][\lforall[x][(x \in y \liff x \notin x)]],
\]
یعنی وجود مجموعهٔ همهٔ مجموعه‌هایی را بیان می‌کند که !!{element}s
خودشان نیستند. چنین مجموعه‌ای نمی‌تواند وجود داشته باشد---این همان
پارادوکس راسل است. در واقع، $\Log{ZFC}$ نسخه‌ای مقید---و سازگار---از
این اصل، یعنی اصل جداسازی، را در بر دارد:
\[
\lforall[z][\lexists[y][\lforall[x][(x \in y \liff (x \in z \land
      !A(x)))]]].
\]

\begin{prob}
با ارائهٔ !!a{derivation} که رابطهٔ زیر را نشان می‌دهد، اثبات کنید اصل
مجموعه‌سازی ناسازگار است:
\[
\lexists[y][\lforall[x][(x \in y \liff x \notin x)]] \Proves \lfalse.
\]
ممکن است مفید باشد نخست نشان دهید
$(A \lif \lnot A) \land (\lnot A \lif A) \Proves \lfalse$.
\end{prob}
\end{document}
